\begin{table}[H]
\centering
\caption{Sensitivity to Auxiliary Scale Regularization}
\label{tab:struct_regularization_sensitivity}
\begin{threeparttable}
\footnotesize
\begin{tabular}[t]{lccccc}
\toprule
Specification & Raw norm & $\lambda$ prior & $\varphi$ prior & Max ITT/SE & High-input \\
\midrule
Preferred & 0.05 & 10.0 & 2.0 & 0.281 & $+0.187$ \\
No raw-norm penalty & 0.00 & 10.0 & 2.0 & 0.292 & $+0.191$ \\
Weak scale priors & 0.05 & 2.5 & 0.5 & 0.273 & $+0.180$ \\
No lambda/phi priors & 0.05 & 0.0 & 0.0 & 0.283 & $+0.242$ \\
No auxiliary scale discipline & 0.00 & 0.0 & 0.0 & 0.289 & $+0.068$ \\
\bottomrule
\end{tabular}
\begin{tablenotes}[para,flushleft]
\footnotesize
\item \textit{Notes:} Each row re-estimates the stage-4 production block holding the residual prior at its preferred value and varying only auxiliary scale regularization. ``Raw norm'' is the coefficient on the raw-parameter norm penalty; ``$\lambda$ prior'' and ``$\varphi$ prior'' are the weights on scale priors centered at 0.4 and 0.2, respectively. The high-input ATE sets $\rho$ to the highest observed signal quality, $\omega=0.98$, and $\tau=0.95$.
\end{tablenotes}
\end{threeparttable}
\end{table}
